% Part: lambda-calculus
% Chapter: introduction
% Section: fixed-point-combinator

\documentclass[../../../include/open-logic-section]{subfiles}

\begin{document}

\olfileid{lam}{int}{fix}
\olsection{स्थिरबिंदू संयोजक}

तुमच्याकडे लॅम्डा पद $g$ आहे आणि दुसरे पद $k$ हवे
आहे असे समजा; त्यात $k$ हे $gk$ शी
$\beta$-सममूल्य असावे.
%READERNOTE{SOL6-A113: दिलेल्या पदातील मुक्त चर पकडले जाऊ नयेत म्हणून diagonal construction मधील अमूर्त चर आधी ताजे निवडले आहे. सर्व fixed-point expressions जपली आहेत.}
आपल्या संकेतनरूढींनुसार पुढील पदे परिभाषित करा;
येथील अमूर्त केलेले चर दिलेल्या पदात मुक्त नसावे.
आवश्यक असल्यास त्याचे आधी नामांतर करा:
\[
\fn{diag}(x) = xx
\]
आणि
\[
l(x) = g(\fn{diag}(x))
\]
म्हणजे $l$ हे पद $\lambd[x][g(xx)]$ आहे.
$k$ हे पद $ll$ असू द्या. मग आपल्याला मिळते:
\begin{align*}
k & = (\lambd[x][g(xx)])(\lambd[x][g(xx)]) \\
& \red  g((\lambd[x][g(xx)])(\lambd[x][g(xx)])) \\
& = gk.
\end{align*}
आता
\[
Y = \lambd[g][((\lambd[x][g(xx)])(\lambd[x][g(xx)]))]
\]
असे घेतल्यास $Yg$ आणि $g(Yg)$ यांचे न्यूनीकरण
होऊन एकच पद मिळते; म्हणून $Yg \equiv_\beta
g(Yg)$. याला “करीचा संयोजक” म्हणतात. त्याऐवजी
\[
Y = (\lambd[xg][g(xxg)])(\lambd[xg][g(xxg)])
\]
असे घेतल्यास प्रत्यक्षात $Yg$ चे $g(Yg)$ पर्यंत
न्यूनीकरण होते; हे अधिक प्रबळ विधान आहे.
$Y$ च्या या दुसऱ्या रूपाला “ट्यूरिंगचा संयोजक”
म्हणतात.

\end{document}
